% Equations of the tqtoy model. Source of docs/MODEL_EQUATIONS.pdf.
% Build:  pdflatex model_equations.tex
\documentclass[11pt,a4paper]{article}
\usepackage[T1]{fontenc}
\usepackage[utf8]{inputenc}
\usepackage{amsmath,amssymb}
\usepackage[margin=2.5cm]{geometry}
\usepackage{booktabs}
\usepackage{array}
\usepackage{url}

% Flat layout: no paragraph indentation, blank line between paragraphs.
\setlength{\parindent}{0pt}
\setlength{\parskip}{0.8em}
\linespread{1.05}

% \code{...} writes a configuration path verbatim and lets it break after '.' and '_',
% which a plain \texttt cannot do and which produces overfull lines.
\DeclareUrlCommand\code{\urlstyle{tt}}
\newcommand{\dd}[2]{\frac{\mathrm{d}#1}{\mathrm{d}#2}}
\newcommand{\h}{\mathrm{h}}
\newcommand{\co}{\mathrm{c}}

\title{\texttt{tqtoy}: model equations}
\date{}

\begin{document}
\maketitle

Everything \texttt{tqtoy} solves is on the following pages. One idea per paragraph, one
equation per box. Configuration parameters are written in \texttt{typewriter}.

\section{What the model is}

Zero-dimensional. Every quantity is a local density, a temperature or a power per unit volume.
Temperatures are in eV, densities in $\mathrm{m^{-3}}$, powers in $\mathrm{W\,m^{-3}}$.

A thermal energy density is $\tfrac{3}{2} e\, n T$.

The plasma has one or two thermal populations (\code{model.populations}). With two, the
pre-disruption plasma is the \textbf{hot} one, subscript $\h$, and the material deposited by the
injection forms the \textbf{cold} one, subscript $\co$. Both occupy the same volume. With one
population everything lives in the cold slots.

Throughout, $\alpha \equiv 2/(3e)$ converts a power density into a $\mathrm{d}(nT)/\mathrm{d}t$.

\section{Symbols}

\begin{center}
\begin{tabular}{@{}l>{\raggedright\arraybackslash}p{9.2cm}l@{}}
\toprule
Symbol & Meaning & Unit \\
\midrule
$n_{e,p}$, $n_{i,p}$ & electron and main-ion density of population $p$ & $\mathrm{m^{-3}}$ \\
$T_{e,p}$, $T_{i,p}$ & electron and main-ion temperature of population $p$ & eV \\
$n_{k,j}$ & density of charge state $k$ of impurity element $j$ & $\mathrm{m^{-3}}$ \\
$Z_j$ & atomic number of element $j$ & \\
$S_{k,j}$, $a_{k,j}$ & ADAS ionisation (\textsc{scd}) and recombination (\textsc{acd}) rates & $\mathrm{m^{3}\,s^{-1}}$ \\
$L^{\mathrm{line}}$, $L^{\mathrm{cont}}$ & ADAS line (\textsc{plt}) and continuum (\textsc{prb}) rates & $\mathrm{W\,m^{3}}$ \\
$S_D$ & deuterium deposition rate & $\mathrm{m^{-3}\,s^{-1}}$ \\
$S^{\mathrm{imp}}_j$ & impurity deposition rate of element $j$ & $\mathrm{m^{-3}\,s^{-1}}$ \\
$J$, $E$ & current density, parallel electric field & $\mathrm{A\,m^{-2}}$, $\mathrm{V\,m^{-1}}$ \\
$\eta$, $\sigma$ & resistivity, conductivity & $\Omega\,\mathrm{m}$, $\mathrm{S\,m^{-1}}$ \\
$\tau_{\alpha\beta}$ & energy-exchange time between species $\alpha$ and $\beta$ & s \\
$e$ & elementary charge, $1.602\times 10^{-19}$ & C \\
\bottomrule
\end{tabular}
\end{center}

\section{State vector}

The solver integrates
\[
  y = \bigl(\, T_{e,\h},\ T_{e,\co},\ n_{e,\h},\ n_{e,\co},\
               T_{i,\h},\ T_{i,\co},\ n_{i,\h},\ n_{i,\co},\
               n_{0,0},\ \dots,\ n_{Z_0,0},\ n_{0,1},\ \dots \,\bigr).
\]

Eight fluid quantities, then every charge state of every impurity element.

Internally $y$ is normalised by fixed scales: $10^3$ for temperatures, $10^{20}$ for fluid
densities, $10^{19}$ for charge states. This is bookkeeping only and does not change the
equations.

\section{Impurity charge states}

Non-coronal balance between electron-impact ionisation and recombination:
\begin{equation}
\boxed{\;
  \dd{n_{k,j}}{t} = r\, n_e \left[\, S_{k-1,j}\, n_{k-1,j} - S_{k,j}\, n_{k,j}
                  + a_{k+1,j}\, n_{k+1,j} - a_{k,j}\, n_{k,j} \,\right]
                  + \delta_{k,q}\, S^{\mathrm{imp}}_j \;}
\end{equation}

Out-of-range coefficients are zero: $S_{-1,j} = S_{Z_j,j} = a_{0,j} = a_{Z_j+1,j} = 0$.

$r = $ \code{model.atomic_rate_multiplier} scales both rates. Large values push the
distribution towards coronal equilibrium.

$q = $ \code{injection.charge_state} is the charge state in which the injected impurities
arrive, $0$ by default.

With two populations each one drives its own ionisation and recombination, and the two
contributions add. The impurity source is counted once.

The inventory of each element, $\sum_k n_{k,j}$, is constant outside the injection window.

\section{Particle balances}

Every injected deuteron, and every electron released by impurity ionisation, is added to the
\textbf{cold} population. No particle is transferred between the two populations, so the hot
one is closed:
\begin{equation}
\boxed{\;
  \dd{n_{e,\co}}{t} = S_D + \sum_{k,j} k\, \dd{n_{k,j}}{t} ,
  \qquad
  \dd{n_{i,\co}}{t} = S_D ,
  \qquad
  \dd{n_{e,\h}}{t} = \dd{n_{i,\h}}{t} = 0 . \;}
\end{equation}

The sum counts the free electrons released, or captured, by the charge-state balance.

Charge neutrality is an exact invariant of these equations:
\begin{equation}
  n_{e,\h} + n_{e,\co} \;=\; n_{i,\h} + n_{i,\co} + \sum_{k,j} k\, n_{k,j}
  \qquad \text{at all times.}
\end{equation}

The initial state satisfies it, so any departure signals a bug rather than a physical effect.

\section{Energy balances}

For each population $p$ and each species:
\begin{equation}
\boxed{\;
  \dd{\left(n_{e,p} T_{e,p}\right)}{t}
    = \alpha \left[\, P^{\mathrm{ohm}}_p - P^{\mathrm{rad}}_p - P^{\mathrm{stoch}}_p
                    + P^{\mathrm{exch}}_{e,p} - c\, T_{e,p} \,\right] \;}
\end{equation}
\begin{equation}
\boxed{\;
  \dd{\left(n_{i,p} T_{i,p}\right)}{t}
    = \alpha \left[\, P^{\mathrm{exch}}_{i,p} - c\, T_{i,p} \,\right] \;}
\end{equation}

$c = $ \code{model.linear_loss_coefficient} is an optional loss proportional to the
temperature, off by default.

The four power terms are defined in the next section.

\textbf{Two ways to integrate the same thing.} \code{solver.state_variables} chooses what
the four thermal slots hold. With \texttt{energy} they hold $nT$ and the equations above are
integrated as written. With \texttt{temperature}, the default, they hold $T$ and the same
equations are divided by $n$:
\begin{equation}
  \dd{T_{e,p}}{t} = \frac{1}{n_{e,p}}
    \left\{ \alpha \left[\, \dots \,\right] - T_{e,p}\, \dd{n_{e,p}}{t} \right\} .
\end{equation}

The last term is the dilution by the injected material, which arrives with no thermal energy.

The energy form has no $1/n$ and is the regular one. It is also the harder one to integrate,
because $nT$ spans about fourteen decades over a run while $T$ spans two, and a single absolute
tolerance cannot cover the former. Hence the default. Outputs are temperatures either way.

\section{Power terms}

\subsection*{Ohmic heating}

\begin{equation}
  P^{\mathrm{ohm}}_p = \eta_p J_p^2 = E\, J_p .
\end{equation}

$J$ is constant by default. \code{plasma.J_final} and \code{plasma.J_ramp_time} add an
optional linear ramp.

\subsection*{Impurity radiation}

Line radiation of the states $0 \dots Z_j - 1$, plus continuum and recombination radiation of
the states $1 \dots Z_j$:
\begin{equation}
  P^{\mathrm{rad}}_p = n_{e,p} \sum_j \left[
      \sum_{k=0}^{Z_j-1} n_{k,j}\, L^{\mathrm{line}}_{k,j}
    + \sum_{k=1}^{Z_j}   n_{k,j}\, L^{\mathrm{cont}}_{k-1,j} \right] .
\end{equation}

Both rates are evaluated at $(n_{e,p}, T_{e,p})$.

Below \code{model.radiation_limit.Te_below} $= 100$~eV the radiated power is capped at the
electron thermal energy divided by \code{model.radiation_limit.time_scale} $= 1\,\mu$s:
\begin{equation}
  P^{\mathrm{rad}}_p \leftarrow \min\!\left( P^{\mathrm{rad}}_p ,\;
    \tfrac{3}{2}\, n_{e,p} T_{e,p}\, e / \tau_{\mathrm{cap}} \right) .
\end{equation}

The limiter can be switched off with \code{model.radiation_limit.enabled}.

\subsection*{Stochastic transport loss}

Electron heat loss along stochastic field lines, a Rechester--Rosenbluth estimate of the
Ward--Wesson type:
\begin{equation}
  P^{\mathrm{stoch}}_p = 2 A\, \frac{n_{e,p}}{r\, a^2}\,
    \frac{1 - 6 (r/a)^2}{\left[1 - (r/a)^2\right]^2}\, \left(T_{e,p}\, e\right)^{7/2} ,
  \qquad
  A = \left(\frac{\delta B}{B}\right)^{2}
      \frac{18\, (2\pi)^{3/2}\, \varepsilon_0^2}{n_{e,p}\, e^4\, \ln\Lambda\, \sqrt{m_e}} .
\end{equation}

The density cancels, so this loss does not depend on $n_{e,p}$.

It assumes $T_e(r) = T_e(0)\,(1 - r^2/a^2)$ and a flat density, and is evaluated at the minor
radius $r$. The profile model requires $r/a < 1/\sqrt{6}$, which the configuration checks. The
Coulomb logarithm here is the constant \code{model.stochastic.coulomb_log} $= 15$.

\subsection*{Collisional exchange}

Energy-exchange time between species $\alpha$ and $\beta$:
\begin{equation}
  \tau_{\alpha\beta} = \frac{3\sqrt{2}\, \pi^{3/2}\, \varepsilon_0^2\, m_\alpha m_\beta}
                            {n_\beta\, e^4\, \ln\Lambda}
    \left( \frac{T_\alpha e}{m_\alpha} + \frac{T_\beta e}{m_\beta} \right)^{3/2} .
\end{equation}

Power transferred from $\alpha$ to $\beta$:
\begin{equation}
\boxed{\;
  P_{\alpha \to \beta} = \frac{3\, n_\alpha\, e\, \left(T_\alpha - T_\beta\right)}
                              {2\, \tau_{\alpha\beta}} \;}
\end{equation}

Positive when $\alpha$ is the hotter, and zero when the two temperatures are equal. That second
property is what makes the two-population coupling of the next section a relaxation.

The Coulomb logarithms follow the usual formularies, each floored at 1 and evaluated with the
hotter population:
\begin{align}
  \ln\Lambda_{ee} &= \begin{cases}
      23.0000 - \ln\!\left( \sqrt{n_e^{\mathrm{cgs}}}\ T_e^{-3/2} \right), & T_e < 10\ \mathrm{eV}, \\
      24.1513 - \ln\!\left( \sqrt{n_e^{\mathrm{cgs}}}\ T_e^{-1} \right),   & T_e \geq 10\ \mathrm{eV},
    \end{cases} \\
  \ln\Lambda_{ei} &= 15.2 - \tfrac{1}{2}\ln\!\left(n_e / 10^{20}\right) + \ln\!\left(T_e / 10^{3}\right), \\
  \ln\Lambda_{ii} &= 17.3 - \tfrac{1}{2}\ln\!\left(n_i / 10^{20}\right) + \tfrac{3}{2}\ln\!\left(T_i / 10^{3}\right),
\end{align}
with $n_e^{\mathrm{cgs}} = n_e \times 10^{-6}$.

With two populations, six exchange powers are computed. Writing
$P_{ee} = P_{e,\h \to e,\co}$ and $P_{ii} = P_{i,\h \to i,\co}$, the net power received by each
species is
\begin{align}
  P^{\mathrm{exch}}_{e,\h}  &= -P_{ee} - P_{e\h \to i\h} - P_{e\h \to i\co} , \\
  P^{\mathrm{exch}}_{e,\co} &= +P_{ee} - P_{e\co \to i\co} - P_{e\co \to i\h} , \\
  P^{\mathrm{exch}}_{i,\h}  &= -P_{ii} + P_{e\h \to i\h} + P_{e\co \to i\h} , \\
  P^{\mathrm{exch}}_{i,\co} &= +P_{ii} + P_{e\co \to i\co} + P_{e\h \to i\co} .
\end{align}

With one population only $P_{e \to i}$ survives.

\code{model.block_electron_ion_exchange} removes every electron-ion term and keeps
$P_{ee}$ and $P_{ii}$.

\section{Resistivity and current sharing}

Effective charge, counting the main ions and every impurity charge state:
\begin{equation}
  Z_{\mathrm{eff}} = \frac{n_i + \sum_{k,j} k^2\, n_{k,j}}{n_i + \sum_{k,j} k\, n_{k,j}} ,
  \qquad n_i = n_{i,\h} + n_{i,\co} .
\end{equation}

Spitzer resistivity, with a $Z_{\mathrm{eff}}$ correction normalised so that $f(1) = 1$:
\begin{equation}
  \eta(T_e) = 1.65\times 10^{-9}\, \ln\Lambda_{ee}\, \left(T_e / 10^{3}\right)^{-3/2}
              f\!\left(Z_{\mathrm{eff}}\right) ,
  \qquad
  f(Z) = \frac{Z\left(1 + 1.198 Z + 0.222 Z^2\right)}{1 + 2.966 Z + 0.753 Z^2}\,
         \frac{3.719}{1.420} .
\end{equation}

With two populations the current is shared between them. Each carries it in proportion to its
own electron density:
\begin{equation}
\boxed{\;
  \sigma_p = \frac{n_{e,p}}{n_{e,\h} + n_{e,\co}}\, \frac{1}{\eta(T_{e,p})} ,
  \qquad
  J_p = J\, \frac{\sigma_p}{\sigma_\h + \sigma_\co} \;}
\end{equation}

The two conductivities add up to $1/\eta$ when the two temperatures are equal.

The density weight matters. The Spitzer resistivity does not depend on $n_e$, because the $n_e$
of $\sigma = n_e e^2 \tau / m_e$ cancels against $\tau \propto 1/n_e$. That holds for a complete
plasma. For a sub-population the collision frequency is set by the \emph{total} ion density, so
its conductivity is proportional to its own density. \code{model.current_sharing} = \texttt{spitzer}
drops the weight and leaves two parallel resistors $\eta(T_{e,\h})$ and $\eta(T_{e,\co})$,
which give a nearly empty population a finite share of the current.

An alternative hot-electron conductivity is available
(\code{model.resistivity} = \texttt{multi\_species}). It accounts for collisions of hot electrons on
cold electrons and on both ion populations. It carries its own density factor, so the weighting
above is not applied on top of it. This expression is not validated.

\section{The two populations}

The two populations are co-located. They share the volume, the ion background and the current.
They are coupled by the collisional exchange and by nothing else.

From the energy balance with $P^{\mathrm{exch}}_{e,\h} \supset -P_{ee}$ and
$P^{\mathrm{exch}}_{e,\co} \supset +P_{ee}$, the temperature difference obeys
\begin{equation}
\boxed{\;
  \left.\dd{\left(T_{e,\h} - T_{e,\co}\right)}{t}\right|_{P_{ee}}
    = -\alpha\, P_{ee} \left( \frac{1}{n_{e,\h}} + \frac{1}{n_{e,\co}} \right)
    \; \propto \; -\left(T_{e,\h} - T_{e,\co}\right) . \;}
\end{equation}

A genuine relaxation. Restoring on both sides of the crossing, with no free parameter. Its
fixed point $T_{e,\h} = T_{e,\co}$ is the state where the two components describe one
Maxwellian. The ion side is identical with $P_{ii}$.

A continuing source, such as the ablation, reopens the difference and the exchange closes it
again. There is no threshold and no switch.

\textbf{Why there is no particle transfer.} Electrons of the two components do mix, and a
model could be built that moves particles as well as energy. It is tempting to deduce that
particle rate from the exchange power already computed, as
$S_e = \alpha P_{ee} / T_{e,\h}$. That expression must be rejected. It is a conversion, namely
the number of particles carrying $P_{ee}$ at $\tfrac{3}{2} e T_{e,\h}$ each, and not a rate
law. A physical rate has the form $n/\tau$.

Adding it also destroys the relaxation. With the particles leaving the hot component, the
dilution term it creates,
$-T_{e,\h}(-S_e)/n_{e,\h} = +\alpha P_{ee}/n_{e,\h}$, cancels the exchange term identically,
for any value of $P_{ee}$. Nothing is then left to pull the two electron temperatures together.
With the particles moving the other way the same term is doubled instead. Both count one
collisional process twice, since electron-electron collisions relax the temperatures and
redistribute the particles at the same rate $\nu_{ee}$.

Of the two possible bookkeepings of that single process, only the temperature relaxation keeps
both component temperatures determined. A pure migration drives $n_{e,\h} \to 0$ while
$\mathrm{d}T_{e,\h}/\mathrm{d}t = \alpha [\dots]/n_{e,\h}$ diverges, with a numerator that has
no reason to vanish. The model therefore keeps the exchange alone.

One-population runs are unaffected by any of this.

\section{Injection}

Both sources are active only for
$t \in [\,t_{\mathrm{start}},\ t_{\mathrm{start}} + \Delta t\,]$.

\subsection*{Prescribed deposition (default)}

The deposited densities are reached at a constant rate:
\begin{equation}
  S_D = \frac{n_D}{\Delta t} ,
  \qquad
  S^{\mathrm{imp}}_j = \frac{n^{\mathrm{inj}}_j}{\Delta t} .
\end{equation}

$n_D = $ \code{injection.n_D} is the deposited deuterium atom density.

$n^{\mathrm{inj}}_j = $ \code{impurities.<element>.injected_atom_density} is the deposited
density of atoms of element $j$, summed over all charge states.

\subsection*{Pellet ablation}

Scaling law for mixed $\mathrm{D_2}$/impurity pellets:
\begin{equation}
  G(n_e, T_e) = \frac{C\, \lambda(X)}{f_w (1-X) + X}\;
                \frac{n_e^{1/3}\, T_e^{5/3}}{V}\; \sum_s r_s^{4/3}
                \;\left(\frac{2}{\max(B,2)}\right)^{0.843}
                \tanh\!\left(\frac{t}{t_{\mathrm{on}}}\right) .
\end{equation}

$\lambda(X) = 27.0837 + \tan(1.48709\, X)$, with $X$ the molecular fraction of $\mathrm{D_2}$,
$V$ the deposition volume and $r_s$ the radius of each shard.

$C = 8.12 \times 10^{14}$ and $f_w = 20.183/2.016$ for hydrogen, or
$C = 4.062 \times 10^{14}$ and $f_w = 20.183/4.0282$ for deuterium.

When $X$ is not given it is computed as $X = (n_D/2)/(n_D/2 + n^{\mathrm{inj}})$.

The resulting sources are $S^{\mathrm{imp}} = G(1-X)$ and $S_D = 2 G X$. With two populations
the rate is the sum of the ablation driven by each.

Shard radii are either identical or drawn from the Parks distribution,
$p(r) = r \kappa_p^2 K_0(r \kappa_p)$ with
$\kappa_p = \left(6\pi^2 N_s n_{\mathrm{solid}} / N_{\mathrm{atoms}}\right)^{1/3}$, then
rescaled so that the total shard volume equals the pellet volume.

\section{Floor and stopping}

The eight fluid quantities are floored at \code{solver.floor} $= 0.01$ before the right-hand
side is evaluated. At the floor, only the part of a derivative that would push a quantity
further below its own floor is removed:
\begin{equation}
  \dd{n}{t} \leftarrow \max\!\left(\dd{n}{t},\, 0\right) \ \text{ if } n \leq y_{\mathrm{floor}} ,
  \qquad
  \dd{T}{t} \leftarrow \begin{cases}
    0 & \text{if } n \leq y_{\mathrm{floor}} , \\[2pt]
    \max\!\left(\dd{T}{t},\, 0\right) & \text{if } T \leq y_{\mathrm{floor}} .
  \end{cases}
\end{equation}

A temperature is held by the density floor, because its equation carries a $1/n$ and because in
the energy formulation $T$ is read back as $(nT)/n$, a ratio with no content once $n$ is pinned.

A density is never held by a temperature. A population whose temperature has collapsed
therefore keeps collecting the particles injected into it.

The run stops on a terminal event when a fluid quantity becomes significantly negative. The
exact solution cannot go below its floor, so the event is offset by the absolute tolerance of
the solver, which is the smallest value it resolves.

\section{Diagnostics}

Density-weighted electron temperature:
\begin{equation}
  \langle T_e \rangle = \frac{n_{e,\h} T_{e,\h} + n_{e,\co} T_{e,\co}}{n_{e,\h} + n_{e,\co}} .
\end{equation}

The thermal-quench time $t_{\mathrm{TQ}}$ is the first output time at which
$\langle T_e \rangle$ falls below 100~eV.

Cumulated energies are integrated on the internal steps of the solver, using its dense output,
rather than on the output times.

\section*{Appendix: hot-tail runaway estimator}

A separate module post-processes a run (\texttt{tqtoy fhte}). A relativistic test electron
follows the collisional drag and the parallel field,
\begin{equation}
  \dd{p}{t} = -\frac{e E_\parallel}{m_e c} - \frac{\nu_{\mathrm{drag}}(p)}{p^2} ,
\end{equation}
and the trajectory ending at the critical momentum
$p_c = \left(E_\parallel / E_c - 1\right)^{-1/2}$ is integrated backward to the initial time.

The runaway fraction is the part of the initial Maxwell--Jüttner distribution above the initial
momentum of that trajectory. The critical field is the Connor--Hastie value,
\begin{equation}
  E_c = \frac{n_e\, e^3 \ln\Lambda_c}{4\pi \varepsilon_0^2 m_e c^2} ,
  \qquad
  \ln\Lambda_c = 14.6 + \tfrac{1}{2} \ln\!\left( T_e / \left(n_e/10^{20}\right) \right) .
\end{equation}

The options are documented in the docstring of \code{src/tqtoy/fhte.py}.

\end{document}
